% !TEX root = main.tex
% Общие поля задания и текста ВКР. msp в классе: 01.04.02, профиль «Системное программирование».
\year{2026-2027}

\student{Синев Денис Евгеньевич
\brief Д.\,Е.\,Синев
\group М-25
\english Sinev Denis Evgenievich
}

\superviser{Секорин Всеслав Станиславович
\brief В.\,С.\,Секорин
\position доцент
\degree к.\,ф.-м.\,н.
}

\program{msp}

\title{Разработка системы мониторинга микросервисов с элементами машинного обучения
    \english Development of a microservice monitoring system with machine learning elements}

\goal{Разработка системы, которая приводит телеметрию микросервисов к единому виду, обнаруживает аномалии алгоритмом без учителя и может работать рядом с Prometheus и Zabbix.}

\begin{tasks}
\task{Реализовать ETL-конвейер для бенчмарков NAB и KPI и для промышленной телеметрии Prometheus и Zabbix: единая схема и одно признаковое пространство.
\result{Построен конвейер на Python~3.12. NAB, KPI, Prometheus и Zabbix приводятся к шести колонкам и к одним и тем же признакам ряда. Исправленный unified-прогон содержит 6\,292\,551 строку.}}
\task{Сравнить алгоритмы обнаружения аномалий без учителя на размеченных рядах при временном сплите и выбрать модель.
\result{Двенадцать детекторов сравнены на 58 рядах NAB и на восьми рядах KPI. По PR-AUC выбрана Isolation Forest: 0{,}207 на NAB и 0{,}464 на срезе KPI.}}
\task{Обучить выбранную модель и сохранить воспроизводимые артефакты.
\result{Сводка, кривые и матрицы ошибок записаны в \texttt{reports/ml}. Детектор учится отдельно на каждом ряде и не видит метку.}}
\task{Оформить сервис с конфигурируемым источником метрик и каналом оповещений.
\result{Зафиксированы контракт конфигурации, то же признаковое окно и канал webhook. Рядом уже работают конвейер \texttt{main.py} и программа сравнения \texttt{python -m ml compare}.}}
\end{tasks}
